\chapter{Аналитическая часть}\label{ch:-4}

В данном разделе будут рассмотрены теоретические аспекты алгоритмов умножения матриц.


\section{Стандартный алгоритм}

Стандартный алгоритм умножения матриц \( A \) и \( B \) предполагает, что для матриц размеров \( m \times n \) и \( n \times p \) результирующая матрица \( C \) будет иметь размер \( m \times p \). Элементы результирующей матрицы \( C \) вычисляются по следующей формуле:

\begin{equation}
    c_{ij} = \sum_{k=1}^{n} a_{ik} \cdot b_{kj},
\end{equation}

где \( c_{ij} \) — элемент матрицы \( C \), полученный умножением \( i \)-й строки матрицы \( A \) и \( j \)-го столбца матрицы \( B \).


\section{Алгоритм Винограда}

Алгоритм Винограда является оптимизированным методом умножения матриц, который уменьшает количество операций умножения за счёт предварительных вычислений. В этом алгоритме для ускорения используются дополнительные массивы \( \text{row\_factor} \) и \( \text{col\_factor} \).

Пусть \( A \) — матрица размером \( m \times n \), \( B \) — матрица размером \( n \times p \), и \( C \) — результирующая матрица размером \( m \times p \).

Рассчитываются вспомогательные массивы:

\begin{equation}
    \text{row\_factor}_i = \sum_{k=1}^{\lfloor n/2 \rfloor} a_{i, 2k - 2} \cdot a_{i, 2k - 1},
\end{equation}

\begin{equation}
    \text{col\_factor}_j = \sum_{k=1}^{\lfloor n/2 \rfloor} b_{2k - 2, j} \cdot b_{2k - 1, j}.
\end{equation}

Элементы результирующей матрицы \( C \) вычисляются следующим образом:

\begin{equation}
    \begin{split}
        c_{ij} = -\text{row\_factor}_i - \text{col\_factor}_j \\
        + \sum_{k=1}^{\lfloor n/2 \rfloor} \Big( (a_{i, 2k - 2} + b_{2k - 1, j})
        \cdot (a_{i, 2k - 1} + b_{2k - 2, j}) \Big).
    \end{split}
\end{equation}

Если \( n \) нечётное, необходимо добавить корректирующий член для учёта последнего элемента строки и столбца:

\begin{equation}
    c_{ij} += a_{i, n-1} \cdot b_{n-1, j}.
\end{equation}

\section*{Вывод}

В данном разделе были рассмотрены теоретические аспекты алгоритмов умножения матриц.

